\selectlanguage{english}

\sheettitle
  {Séance of 14 July 2026}
  {$L^p$ spaces, density measures, and Radon--Nikodym}

This is not a construction of the Radon--Nikodym theorem. The theorem is treated as prior knowledge and recalled only at the level of hypotheses, reference measure, signed extension, and almost-everywhere uniqueness. The work of this session is to recover the $L^p$ structure used later for expectation, conditioning, conditional expectation, and variance.

\section*{Order of ideas}

\[
L^p(\mu)\text{ as an ambient-measure-dependent object}
\longrightarrow
\text{probability-space moment control}
\longrightarrow
\text{bounded tests}
\longrightarrow
\text{density-generated measures}
\longrightarrow
\text{Radon--Nikodym representations}.
\]

There is no implication
\[
L^q(\mu)\subseteq L^p(\mu)
\Longrightarrow
\text{Hahn local domination}.
\]
The two statements belong to different proofs. The Hahn-first construction remains part of later study.

\section*{Programme: 90 minutes}

\begin{enumerate}
  \item \textbf{15 min.} Complete Exercise 0: recover the ambient measure, quotient relation, moment interpretation, and the roles of $L^1$, $L^\infty$, and $L^2$.
  \item \textbf{20 min.} Complete Exercise 1: prove the finite-measure embedding and its $q=\infty$ endpoint. Retain the Hölder exponents and the exact norm constant.
  \item \textbf{20 min.} Complete Exercise 2: pass from indicator tests to bounded measurable tests, naming dominated convergence at the final passage.
  \item \textbf{20 min.} Complete Exercise 3: construct a finite signed measure from an $L^1$ density, prove density uniqueness, and state change of measure in its legitimate integrability regime.
  \item \textbf{10 min.} Complete Exercise 4: identify the numerator, reference measure, measurable space, and almost-everywhere basis in four later uses of the theorem.
  \item \textbf{5 min.} Record the first rupture and assign the session verdict.
\end{enumerate}

\section*{Required reconstruction}

On a probability space,
\[
L^\infty(\mathbb P)
\subseteq
L^2(\mathbb P)
\subseteq
L^1(\mathbb P),
\]
and
\[
L^1\times L^\infty\longrightarrow L^1,
\qquad
L^2\times L^2\longrightarrow L^1.
\]
For a general measure space,
\[
f\in L^1(\mu)
\longmapsto
\nu_f=f\mu
\]
produces a finite signed measure, while
\[
\nu=h\mu
\Longrightarrow
\int \varphi\,d\nu=\int \varphi h\,d\mu
\]
holds for nonnegative $\varphi$ and, in the signed regime, under the corresponding absolute-integrability hypothesis.

\section*{Radon--Nikodym review}

The session requires only the following distinctions.
\begin{enumerate}
  \item Forward density generation and the RN converse run in opposite directions.
  \item In the standard positive theorem, the numerator and reference measures are $\sigma$-finite and the numerator is absolutely continuous with respect to the reference measure.
  \item A finite signed numerator is handled through its Jordan parts.
  \item The derivative is unique only relative to the reference measure, almost everywhere.
\end{enumerate}

\section*{Further work}

Sharpness of the embedding constant, long infinite-measure counterexamples, the full Young--Hölder--Minkowski reconstruction, the Hahn local-density fragment, maximal represented mass, residue elimination, $\sigma$-finite gluing, and full $L^2$ projection theory remain for the extended course of study.

Solutions are forbidden before the first complete written attempt. A familiar theorem name with no numerator or reference measure attached is not yet a representation.
